\chapter{Связывающая форма Дирихле и квантово-марковская полугруппа}
\label{ch:v8-linking-dirichlet-quantum-markov-semigroup-gate}

\section{Алгебра наблюдаемых уже присутствует}

Предыдущий гейт показал, что полное корреляционное ядро имеет смысл лишь
вместе с алгеброй различимых наблюдаемых. В связывающем секторе Тома VII
уже имеется каноническая концевая декомпозиция
\begin{equation}
 \mathcal A_{\rm end}=M_{11}(\mathbb C)\oplus M_{10}(\mathbb C),
 \label{eq:v8-endpoint-observable-algebra}
\end{equation}
соответствующая исходных  и целевых модулям физической инцидентности
$A_0:\mathbb C^{11}\to\mathbb C^{10}$. Это не коммутативная координатная
алгебра и не выбранный диагональный базис, но это честная алгебра
наблюдаемых концевых секторов, возникающая как два угла связывающего носителя.

Определим самосопряжённый связывающий оператор
\begin{equation}
 D_A=\begin{pmatrix}0&A_0^*\\A_0&0\end{pmatrix}
 \quad\text{на}\quad \mathbb C^{11}\oplus\mathbb C^{10}.
 \label{eq:v8-linking-dirac}
\end{equation}

\section{Каноническая диффузия наблюдаемых}

Квадратичная форма коммутатора задаёт генератор
\begin{equation}
 \mathcal L_A(Z)=-\frac12[D_A,[D_A,Z]].
 \label{eq:v8-linking-qms-generator}
\end{equation}
Подалгебра~\eqref{eq:v8-endpoint-observable-algebra} инвариантна. Для
$Z=X\oplus Y$ ограничение имеет явный вид
\begin{equation}
 \mathcal L_A(X,Y)=
 \left(
 A_0^*YA_0-\frac12\{A_0^*A_0,X\},
 A_0XA_0^*-\frac12\{A_0A_0^*,Y\}
 \right).
 \label{eq:v8-linking-corner-markov-generator}
\end{equation}
Перекрёстные слагаемые действительно переносят наблюдаемое из одного угла
в другой; это не тензорное продолжение независимых секторных полугрупп.

Генератор имеет стандартную форму с одним самосопряжённым Линдблад-
оператором $D_A$. Поэтому
\begin{equation}
 T_t=e^{t\mathcal L_A},\qquad t\ge0,
 \end{equation}
является вполне положительной унитальной полугруппой, сохраняющей полный
след. В собственном базисе $D_A$ она умножает матричную единицу $E_{ij}$ на
\begin{equation}
 \exp\left[-\frac t2(\lambda_i-\lambda_j)^2\right],
 \label{eq:v8-linking-gaussian-schur-kernel}
\end{equation}
а матрица этих коэффициентов положительна как гауссово ядро. Машинный
аудит независимо проверяет унитальность, сохранение следа, полугрупповое
тождество, отрицательность генератора и ковариантность при независимых
заменах исходных  и целевых базисов.

\section{Неподвижная алгебра}

Положительный результат не устраняет кратности. Сингулярный спектр $A_0$
содержит значение $1$ с кратностью шесть, четыре простых ненулевых значения
и одномерное исходное ядро. Поэтому размерность неподвижного пространства
равна
\begin{equation}
 \dim\ker\mathcal L_A
 =6^2+4\cdot1^2+1^2=41.
 \label{eq:v8-linking-fixed-algebra-dimension}
\end{equation}
Прямая диагонализация матрицы генератора размера $221$ подтверждает это
точно в численной точности. Щель над неподвижной алгеброй положительна и
равна $0.01398359846\ldots$, но сама полугруппа не примитивна и не имеет
единственного стационарного наблюдаемого.

\begin{theorem}[Положительный допуск квантовой диффузии]
Физическая инцидентность Тома VII без новых коэффициентов задаёт на
$M_{11}\oplus M_{10}$ симметричную унитальную и сохраняющую след
квантово-марковскую полугруппу с ненулевым двусторонним переносом между
концевыми углами. Однако её неподвижная алгебра имеет размерность $41$.
Следовательно, динамика переходов получена, а коммутативная алгебра событий
и локальный селектор ещё не выведены.
\end{theorem}

\begin{nogocriterion}[Запрет выбора MASA после результата]
Нельзя после вычисления выбрать удобную максимальную коммутативную
подалгебру в 41-мерном неподвижном секторе и назвать её пространством
событий. Допустимое продолжение обязано показать, что уже выведенные калибровочные,
Ходж или аффинные операторы канонически снимают шестикратное вырождение и
сохраняют вполне положительную полугруппу.
\end{nogocriterion}

\begin{protocol}
\begin{enumerate}
 \item концевая алгебра $M_{11}\oplus M_{10}$: \textbf{имеется};
 \item минимальные диагональные проекторы назначены: \textbf{нет};
 \item генератор $-\frac12\operatorname{ad}_{D_A}^2$: \textbf{выведен};
 \item полная положительность и унитальность: \textbf{пройдены};
 \item сохранение следа: \textbf{пройдено};
 \item двусторонний перенос между углами: \textbf{ненулевой};
 \item размерность неподвижной алгебры: \textbf{$41$};
 \item единственное равновесие: \textbf{нет};
 \item коммутативная локальность: \textbf{не выведена};
 \item квантово-марковский маршрут: \textbf{положительно открыт}.
\end{enumerate}
\end{protocol}

Машинный сертификат сохранён в
\path{s2t_v8_linking_dirichlet_quantum_markov_semigroup_gate_results.json}.

\begin{thebibliography}{9}
\bibitem{v8-qms-generators}
G.~Androulakis, M.~Ziemke,
\emph{Generators of Quantum Markov Semigroups}, arXiv:1406.3417.
\bibitem{v8-qms-deschamps}
J.~Deschamps, F.~Fagnola, E.~Sasso, V.~Umanit\`a,
\emph{Structure of uniformly continuous quantum Markov semigroups},
arXiv:1412.3239.
\bibitem{v8-qms-wirth}
M.~Wirth,
\emph{A noncommutative transport metric and symmetric quantum Markov
semigroups as gradient flows of the entropy}, arXiv:1808.05419.
\end{thebibliography}